\chapter{The Intent, the Contract and the Export}
\label{ch:export}

\section{What the Intent States}

The intent is the third of the four documents: teams, spawns, protections, build
regions, objectives, and how each objective is captured. It is the only one of the
four that is about the \emph{match} rather than about the ground.

Storing it and projecting it into the PGM document is one call and one idempotent
pass. \route{PUT /api/map/\{slug\}/intent} writes the intent and derives from it the
teams, kits, regions, filters, apply-rules and spawns that the document needs.
\Cref{ch:division} gave the scale of that derivation: a plan stating four wool rooms
and their markers produced four wools, four monuments, 22 regions, 34 filters and 11
apply-rules, with nothing authored and nothing corrected.

\section{Naming, and a Trap in It}

PGM has no \id{<monument>} element. The destroy objective's element is
\id{<destroyable>}, and the word \emph{monument} is already taken throughout the
studio for the CTW wool-monument, which is the block a capturing team places wool
on.

\begin{ruling}{The Objective Is a Destroyable, Never a Monument}
Colloquially a destroyable \emph{is} the DTM monument. In types, columns,
identifiers and JSON keys it is not, because the word already names a different
thing in the same system.
\end{ruling}

This paper uses \emph{monument} in the colloquial sense when describing how a map
plays, because that is what players call it, and \emph{destroyable} whenever the
subject is the document.

\section{The Two Destroy Objectives}

Both are goals owned by one team and destroyed by everyone else, and both are
structurally simpler than a wool, which drags a room region, a spawner, a dye item,
team filters, apply rules and a coupling to the spawns.

\subsection{The Destroyable}

\begin{quote}
\id{<destroyable owner name region materials [completion] [id] [modes]>}. PGM builds
the set of blocks inside the region that match \id{materials}, and the goal
completes when \id{completion} of them are broken.
\end{quote}

\id{completion} defaults to 1.0. Four attributes are required and everything else
defaults.

\subsection{The Core}

\begin{quote}
\id{<core team region [material] [leak] [name] [id] [modes]>}. \id{material}
defaults to obsidian and \id{leak} to 5.
\end{quote}

The core builds two block sets from the world inside its region: the casing, being
the blocks matching \id{material}, and the lava. It leaks when a lava block reaches
a height of \id{region.min.y} minus \id{leak} minus one, within fifteen blocks
horizontally of the core bounds.

One inconsistency is worth noting because it shows in the document: the core's
owning attribute is \id{team} rather than \id{owner}, which is a known
inconsistency in PGM itself. The studio mirrors the XML and calls the field
\id{Owner} in its own code.

\subsection{Why This Caps the Team Count}

\Cref{sec:team-counts} stated the consequence and this is the mechanism. Neither
objective has the wool's per-capturing-team fan-out. The shape is \textbf{one
object, one owner, and $N-1$ attackers}.

A wool at four teams gives each team three distinct things to capture, so each team
has an objective of its own and a winner is well defined. A destroyable at four
teams gives each team one thing to lose and three things to break, and nothing in
that arrangement says who won. It is a property of the goal's structure rather than
a limit anyone imposed.

\section{The Float, and What the Corpus Actually Does}

Both objectives float above the terrain, which is why no carving or negative-space
primitive is needed to seat one. The studio's generator floats a destroyable three
to five blocks, with a default of four, and a lone floating pillar is a normal
output rather than an error. An author's cap bounds the other end at twelve, and the
reasoning behind that cap is that a goal high enough to need a tower under it is
reached by a build project rather than by a raid.

The corpus does not do this, and the difference is worth stating because a detector
expecting a gap will miss most real structures. Measured on the declared blocks
themselves across 571 single-material structures, the median air gap beneath is
\textbf{zero}, and \textbf{424 of 571 rest directly on something}. Only 68 float
three to five blocks. The floating pillar is the authored ideal; the corpus
overwhelmingly seats its destroyables on terrain, a pedestal or a build.

\begin{ruling}{Float Is an Offset, Not a Height}
The studio's \id{float} counts from the ground as built rather than from a
plan-nominal surface. A destroyable or a core is never given a Y at all: the
structure floats its stated offset over whatever the relief actually left under its
column, resolved once the layout is rasterized, which is the only place that ground
is known. A goal placed on a mesa stays on the mesa however the mesa's own height is
later solved.
\end{ruling}

Because the compiler resolves nothing about height for these two markers, they are
also the one marker kind that may name no plan piece at all. The position reads as
an absolute board position, and the ground it rides can be an authored sketch shape
with no piece behind it. A landform is then authored once, as the shape it is,
rather than twice: once as a plan rectangle purely to give a marker something to
ride, and again as the polygon it actually is.

\section{Flooring Is Per Field}

One detail of the contract has cost hours more than once, and it looks like an
inconsistency until the reason is known.

\begin{ruling}{Coordinate Flooring Is Per Field, Not Global}
The wool's \id{<location>} is floored and the destroyable's \id{<block>} is not,
because PGM floors one itself and never the other. Applying one rule globally in
either direction produces a document that is wrong in half its objectives.
\end{ruling}

\section{The Export Gate}

The export is where several checks first become fatal. \route{GET
/api/map/\{slug\}/xml} answers the document, gated on the pre-flight checks, and
\route{GET /api/map/\{slug\}/export} answers the world.

Four refusals are worth knowing by name. \id{EX2} is no spawn of any kind.
\id{EX3} is a kind that went in and did not come out, which is the harder failure
of the two because it means something was accepted and then lost. \id{EX4} is an
objective with no team. \id{OB20}, an unknown gamemode, is checked first on every
export regardless of how the map originated.

Which gate runs does not depend on which door the caller came through. A
sketch-originated map's whole chain runs inside the same composer as a
plan-originated one's.

The corpus records this as an expensive property rather than a neutral one. Two of
those refusals fire only at export, so the first time a board hears about them is a
409 answered \emph{after the whole world has been built}. One run hit them three
times between them, at a cost of three full build cycles.

\section{What Is Delivered}

A sketch-originated map exports a ZIP carrying \id{map.xml}, \id{level.dat} and the
region files at the archive top. A map that arrived through the import path exports
the XML alone.

The world format is Anvil in its numeric block form, which is the 1.8 to 1.12
encoding with nibble-packed sections. That matches the studio's supported range,
and it is the reason the studio refuses maps outside it rather than accommodating
them: a malformed or out-of-range map is rejected rather than handled by weakening
the schema.
